\documentclass[conference]{IEEEtran}

\usepackage{graphicx}
\usepackage{amsmath,amssymb}
\usepackage{float}
\usepackage{booktabs}
\usepackage{multirow}

\begin{document}

\title{
    \textbf{Benchmark Report: 519.lbm\_r} \\
    \text{Register Spill Analysis on RISC-V via gem5}
}

\author{
    \IEEEauthorblockN{Lütfullah Burak Kaya}
    \IEEEauthorblockA{
        Ozyegin University \\
        burak.kaya.50711@ozu.edu.tr
    }
    \and
    \IEEEauthorblockN{Ismail Akturk}
    \IEEEauthorblockA{
        Ozyegin University \\
        ismail.akturk@ozyegin.edu.tr
    }
}

\newcommand{\LongDate}{%
    \number\day\space
    \ifcase\month
    \or January\or February\or March\or April\or May\or June%
    \or July\or August\or September\or October\or November\or December%
    \fi
    \space \number\year
}

\twocolumn[
\begin{@twocolumnfalse}
\maketitle
\begin{center}{\small \LongDate}\end{center}
\vspace{1em}
\end{@twocolumnfalse}
]

% ================================================

\begin{abstract}
This report presents the register spill characterization of the
\texttt{519.lbm\_r} benchmark from SPEC CPU2017, executed on a RISC-V
(RV64GC) target using the gem5 TimingSimpleCPU simulator.
A custom SpillDetector was integrated into gem5 to count store-then-load
patterns on the stack within a user-defined Region of Interest (ROI).
The test dataset (20 time steps, channel flow, $100 \times 100 \times 130$
grid) was used as the input.
Results show a spill rate of \textbf{4.21\%} within the ROI,
corresponding to 274 million detected spill events across 6.5 billion
ROI instructions. The benchmark's lattice Boltzmann fluid solver
iterates over a large 3D grid with 19 distribution functions per cell,
sustaining high register pressure throughout the entire simulation.
\end{abstract}

% =========================================================
\section{Benchmark Overview}

\texttt{519.lbm\_r} is the SPEC CPU2017 rate variant of LBM, a
Lattice Boltzmann Method fluid dynamics simulator. It simulates
incompressible fluid flow in a 3D domain by propagating and
colliding 19 distribution functions per lattice site (D3Q19 model)
using the two-relaxation-time (TRT) collision operator. The benchmark
is dominated by a single innermost kernel ---
\texttt{LBM\_performStreamCollideTRT()} --- which is called once per
time step and sweeps the entire 3D grid.

The RISC-V binary (\texttt{lbm\_r.riscv}) was compiled with gem5 ROI
markers inserted in \texttt{src/main.c}, wrapping the main time-step
loop:

\begin{itemize}
    \item \texttt{m5\_work\_begin(0,0)} --- before the time-step for-loop
    \item \texttt{m5\_work\_end(0,0)}   --- after the time-step for-loop
\end{itemize}

Each iteration of the loop calls \texttt{LBM\_handleInOutFlow()} (for
channel flow simulations) and \texttt{LBM\_performStreamCollideTRT()},
then swaps source and destination grids.

% =========================================================
\section{Simulation Setup}

\subsection{Input Datasets}

The three dataset sizes differ only in the number of time steps
and the initial fluid state (obstacle file), as shown in
Table~\ref{tab:inputs}. The grid dimensions are fixed at
$100 \times 100 \times 130$ cells across all sizes.

\begin{table}[H]
\centering
\caption{lbm\_r Input Dataset Sizes}
\label{tab:inputs}
\begin{tabular}{lcll}
\toprule
\textbf{Dataset} & \textbf{Time steps} & \textbf{Obstacle file} & \textbf{Sim type} \\
\midrule
test    &    20 & 100\_100\_130\_cf\_a.of & channel flow \\
train   &   300 & 100\_100\_130\_cf\_b.of & channel flow \\
refrate & 3,000 & 100\_100\_130\_ldc.of  & lid-driven cavity \\
\bottomrule
\end{tabular}
\end{table}

The \textbf{test} dataset (20 time steps, channel flow) was selected
for this run. The grid has $100 \times 100 \times 130 = 1{,}300{,}000$
lattice sites, each with 19 double-precision distribution values
(19 $\times$ 8 bytes = 152 bytes per site), totalling approximately
198\,MB of working set.

\subsection{gem5 Configuration}

\begin{itemize}
    \item \textbf{CPU model:}   TimingSimpleCPU
    \item \textbf{ISA:}         RISC-V RV64GC
    \item \textbf{Caches:}      enabled (default L1 I/D + L2)
    \item \textbf{Memory:}      4\,GB simulated physical RAM
    \item \textbf{Spill mode:}  summary only (\texttt{spill\_verbose=False})
\end{itemize}

The full gem5 invocation was:

\begin{verbatim}
./build/RISCV/gem5.opt
  --outdir=benchmarks/lbm/m5out_test
  configs/deprecated/example/se.py
  --cmd=benchmarks/lbm/lbm_r.riscv
  --options="20 benchmarks/lbm/reference.dat 0 1
             benchmarks/lbm/100_100_130_cf_a.of"
  --cpu-type=TimingSimpleCPU
  --caches
  --mem-size=4GB
\end{verbatim}

% =========================================================
\section{Simulation Results}

The simulation completed successfully, printing per-step grid
statistics every 64 time steps and producing the reference output file.

\subsection{Pre-ROI Context}

Before the ROI, the binary executed grid allocation, obstacle file
parsing, and fluid initialisation. Table~\ref{tab:preroi} shows
the global counters at ROI entry.

\begin{table}[H]
\centering
\caption{Global Counters at ROI Entry}
\label{tab:preroi}
\begin{tabular}{lr}
\toprule
\textbf{Metric} & \textbf{Value} \\
\midrule
Total instructions (pre-ROI) & 239,085,882 \\
Total spills detected (pre-ROI) & 7,890,132 \\
Pre-ROI spill rate & 3.30\% \\
\bottomrule
\end{tabular}
\end{table}

\subsection{ROI Spill Statistics}

Table~\ref{tab:roi} presents the spill statistics collected within
the ROI (20 LBM time steps).

\begin{table}[H]
\centering
\caption{ROI Spill Statistics --- 519.lbm\_r (test, 20 steps)}
\label{tab:roi}
\begin{tabular}{lr}
\toprule
\textbf{Metric} & \textbf{Value} \\
\midrule
ROI instructions         & 6,506,800,041 \\
ROI stores               &   775,731,184 \\
ROI loads                & 1,276,836,306 \\
\midrule
\textbf{ROI spills}      & \textbf{274,130,440} \\
\textbf{Spill rate}      & \textbf{4.2130\%} \\
\bottomrule
\end{tabular}
\end{table}

\subsection{Timing}

\begin{table}[H]
\centering
\caption{Simulation Timing}
\label{tab:timing}
\begin{tabular}{lr}
\toprule
\textbf{Metric} & \textbf{Value} \\
\midrule
Simulated time           & 37.74\,s \\
Host wall time           & 6,482\,s ($\approx$1.8\,hours) \\
Total simulated instructions & 6,841,278,693 \\
\bottomrule
\end{tabular}
\end{table}

% =========================================================
\section{Analysis}

\subsection{Spill Rate Interpretation}

The measured spill rate of \textbf{4.21\%} means that approximately
1 in every 24 instructions inside the ROI is a register spill.
Table~\ref{tab:compare} places this in context of all evaluated benchmarks.

\begin{table}[H]
\centering
\caption{Cross-Benchmark Spill Rate Comparison}
\label{tab:compare}
\begin{tabular}{lrr}
\toprule
\textbf{Benchmark} & \textbf{ROI Spills} & \textbf{Spill Rate} \\
\midrule
507.cactuBSSN\_r & 3,782,099,438 & 7.8756\% \\
500.perlbench\_r &     1,070,055 & 7.2838\% \\
526.blender\_r   &    61,845,869 & 6.2193\% \\
541.leela\_r     & 1,158,236,271 & 4.4438\% \\
519.lbm\_r       &   274,130,440 & \textbf{4.2130\%} \\
544.nab\_r       &   151,310,616 & 2.2612\% \\
538.imagick\_r   &     2,424,899 & 2.0601\% \\
505.mcf\_r       &   446,113,136 & 1.8263\% \\
508.namd\_r      &   327,718,955 & 1.0269\% \\
\bottomrule
\end{tabular}
\end{table}

\subsection{Store vs.\ Load Balance}

The load count (1.28B) is approximately $1.65\times$ higher than the
store count (776M). In the D3Q19 TRT collision operator, each lattice
site reads distribution values from 19 neighbours (stream step), then
computes the post-collision state and writes it back. The read-heavy
ratio reflects that each of the 19 neighbours is loaded from memory,
while only the local site is written.

\subsection{Spill Fraction of Loads}

Within the ROI, the ratio of spills to total loads is:
\[
\frac{274{,}130{,}440}{1{,}276{,}836{,}306} \approx 21.5\%
\]
Roughly 1 in 5 load instructions is a spill reload. The TRT collision
kernel simultaneously manipulates all 19 distribution values plus
their equilibrium counterparts and collision-relaxation intermediates.
With RISC-V's 32 FP registers unable to accommodate all of these
live values, systematic spilling occurs at every grid point.

\subsection{ROI Coverage}

The ROI accounts for $6{,}506{,}800{,}041$ out of $6{,}841{,}278{,}693$
total simulated instructions:
\[
\frac{6{,}506{,}800{,}041}{6{,}841{,}278{,}693} \approx 95.1\%
\]
The ROI captures almost all execution. The remaining 4.9\% is grid
setup and I/O, confirming the time-step loop dominates the benchmark.

\subsection{Scaling to Refrate}

The refrate dataset uses 3,000 time steps (150$\times$ more than test).
Assuming linear scaling, the estimated refrate spill count is:
\[
274{,}130{,}440 \times 150 \approx 41.1 \text{ billion spills}
\]

% =========================================================
\section{Conclusion}

The \texttt{519.lbm\_r} benchmark exhibits a register spill rate of
4.21\% within its core time-step loop on RISC-V RV64GC. With 274
million spill events across 6.5 billion ROI instructions, the D3Q19
TRT collision kernel sustains continuous register pressure due to the
19 distribution functions and their intermediate values that must be
simultaneously live at each grid point. The ROI captures 95.1\% of
total execution, confirming the time-step loop as the dominant
computation. At 1.8 hours host time for only 20 time steps, a full
refrate run (3,000 steps) would require approximately 270 hours,
making \texttt{519.lbm\_r} one of the most simulation-time-intensive
benchmarks in the SPEC CPU2017 suite under gem5.

\end{document}
